%% ============================================================
%%  Tesis INAOE — documento de ejemplo
%%
%%  Compilar:   latexmk          (lee .latexmkrc: pdflatex + bibtex)
%%  Vigilar:    latexmk -pvc     (recompila al guardar)
%%  Limpiar:    latexmk -C
%%
%%  Compilación rápida mientras se escribe:
%%    \documentclass[12pt,draft]{report} omite imágenes, microtype
%%    e hipervínculos y marca las líneas desbordadas.
%% ============================================================

\documentclass[12pt]{report}

%------------------------------------------------
% Plantilla: fuentes, márgenes, idioma y portada
%------------------------------------------------

\usepackage[english]{inaoe-tesis}
% \usepackage[spanish]{inaoe-tesis}

%------------------------------------------------
% Microtipografía (después de las fuentes)
%------------------------------------------------

\usepackage{microtype}
\setlength{\emergencystretch}{3em}

%------------------------------------------------
% Matemáticas (amsmath ya lo carga la plantilla)
%------------------------------------------------

\usepackage{mathtools}

%------------------------------------------------
% Figuras y tablas
%------------------------------------------------

\usepackage{booktabs}
\usepackage{threeparttable}

%------------------------------------------------
% Algoritmos
%------------------------------------------------

\usepackage{algorithm}
\usepackage{algpseudocode}

%------------------------------------------------
% Código
%------------------------------------------------

\usepackage{listings}

%------------------------------------------------
% Utilidades
%------------------------------------------------

\usepackage{csquotes}
\usepackage[normalem]{ulem}

%------------------------------------------------
% Bibliografía
%------------------------------------------------

% El estilo unsrtnat requiere natbib; numbers -> citas numéricas [1]
\usepackage[numbers]{natbib}

%------------------------------------------------
% Hipervínculos (siempre al final del preámbulo)
%------------------------------------------------

\usepackage[
  colorlinks=true,
  linkcolor=blue,
  citecolor=blue,
  urlcolor=blue,
  hypertexnames=false % evita destinos duplicados al reiniciar la numeración
]{hyperref}

%------------------------------------------------
% Metadatos de la tesis
%------------------------------------------------

\tituloTesis{A Continuous User Authentication Scheme Using PPG Signals from Wearable Devices}

\autor{Diego Cruz Aguilar}

\asesor{
  Dr. Alfonso Martínez Cruz\\
  Dra. Kelsey A. Ramírez Gutiérrez
}

\grado{M.S. in Computer Science}

\mes{October}
\anio{2025}

%================================================
% DOCUMENTO
%================================================

\begin{document}

%------------------------------------------------
% Portada y página en blanco posterior
%------------------------------------------------

\portada
\thispagestyle{empty}\null\clearpage

%------------------------------------------------
% Páginas preliminares
%------------------------------------------------

\pagenumbering{Roman}

\tableofcontents
\listoffigures
\listoftables

%------------------------------------------------
% Dedication
%------------------------------------------------

\chapter*{Dedication}
\addcontentsline{toc}{chapter}{Dedication}

Dedicado a...

%------------------------------------------------
% Acknowledgements
%------------------------------------------------

\chapter*{Acknowledgements}
\addcontentsline{toc}{chapter}{Acknowledgements}

I wish to express my sincere gratitude to...

%------------------------------------------------
% Abstract
%------------------------------------------------

\chapter*{\abstractname}
\addcontentsline{toc}{chapter}{\abstractname}

This thesis presents...

%------------------------------------------------
% Cuerpo principal
%------------------------------------------------

\clearpage
\pagenumbering{arabic}

\chapter{Introduction}

\section{Background}

Text... This document was typeset with \LaTeX~\cite{lamport94}, which is
built on top of \TeX~\cite{knuth84}.

\section{Problem Statement}

Text...

\chapter{Related Work}

Text...

\chapter{Methodology}

Text...

\chapter{Results}

Text...

\chapter{Conclusions and Future Work}

\section{Conclusions}

Text...

\section{Future Work}

Text...

%------------------------------------------------
% Apéndices
%------------------------------------------------

\appendix

\chapter{Published Articles}

\begin{itemize}
  \item \textbf{Article title}, Journal Name, 2025.
\end{itemize}

%------------------------------------------------
% Bibliografía
%------------------------------------------------

\clearpage
\phantomsection % ancla correcta para el enlace del índice
\addcontentsline{toc}{chapter}{\bibname}

\bibliographystyle{unsrtnat}
\bibliography{references}

\end{document}
